% =============================================================================
% Fourier i DETR: wyjaśnienie krok po kroku
% =============================================================================

\section{Od zera do Fouriera: wyjaśnienie krok po kroku}
\label{sec:krok_po_kroku}

Zanim przejdziemy do teorii i wzorów, wyjaśnijmy prostym językiem o~co
w~tym wszystkim chodzi --- od samego początku.


% -----------------------------------------------------------------------------
\subsection{Krok 1: Problem --- za dużo obrazów, za mało czasu}
\label{sec:krok1}

Mamy nagrania wideo z~operacji zaćmy.  Z~tych nagrań wyciągamy klatki
(pojedyncze obrazy).  Po augmentacji (obracanie, przycinanie, zmiana jasności)
powstaje pula \textbf{90\,389 obrazów}.

Problem: trenowanie sieci neuronowej DETR na wszystkich 90k obrazach jest:
\begin{itemize}
    \item \textbf{Wolne} --- DETR trenuje się znacznie dłużej niż YOLO
          (transformery potrzebują więcej czasu na zbieżność).
    \item \textbf{Marnujące zasoby} --- wiele klatek jest prawie identycznych
          (wideo 25--30 kl./s $\Rightarrow$ kolejne klatki różnią się minimalnie).
    \item \textbf{Nieefektywne} --- redundantne dane nie uczą sieci niczego
          nowego, a mogą wręcz pogarszać generalizację.
\end{itemize}

\textbf{Cel:} wybrać \textbf{20\,000 najlepszych obrazów} z~puli 90\,389
(redukcja $4{,}5\times$) tak, żeby DETR trenowany na nich był \emph{lepszy}
niż trenowany na wszystkich.


% -----------------------------------------------------------------------------
\subsection{Krok 2: Czym jest transformata Fouriera na obrazie?}
\label{sec:krok2}

Każdy obraz można rozłożyć na ,,fale'' o różnych częstotliwościach --- tak jak
dźwięk rozkłada się na tony niskie (bas) i wysokie (pisk).

\begin{itemize}
    \item \textbf{Niskie częstotliwości} (pasmo \texttt{low}, 0--10\% promienia)
          = gładkie, wolno zmieniające się obszary obrazu: tło, jednorodna
          tkanka, rozmyte fragmenty.

    \item \textbf{Średnie częstotliwości} (pasmo \texttt{mid}, 10--50\%)
          = tekstury: wzór tęczówki, struktura torebki soczewki, ziarno tkanki.

    \item \textbf{Wysokie częstotliwości} (pasmo \texttt{high}, 50--100\%)
          = ostre krawędzie: brzeg narzędzia chirurgicznego, odbicia światła,
          granice obiektów.
\end{itemize}

Transformata Fouriera (2D FFT) zamienia obraz z~,,pikseli'' na
,,częstotliwości''.  Zamiast pytać ,,jaki kolor ma piksel $(x,y)$?'', pytamy
,,ile energii jest w~falach o~danej częstotliwości i~kierunku?''

\paragraph{Analogia.}
Pomyśl o~equalizerze w~odtwarzaczu muzyki: suwaki bass/mid/treble.
Transformata Fouriera robi dokładnie to samo, ale dla obrazu --- mierzy ile
jest basu (gładkie tło), ile średnich tonów (tekstury) i~ile sopranów (ostre
krawędzie).


% -----------------------------------------------------------------------------
\subsection{Krok 3: Co konkretnie mierzymy? (9 liczb na obraz)}
\label{sec:krok3}

Dla każdego z~90\,389 obrazów obliczamy \textbf{9-wymiarowy wektor cech
Fouriera}:

\begin{table}[h]
\centering
\caption{9 cech Fouriera obliczanych dla każdego obrazu.}
\label{tab:fourier_features_simple}
\begin{tabular}{clp{7.5cm}}
\toprule
\textbf{Nr} & \textbf{Cecha} & \textbf{Co oznacza (prosto)} \\
\midrule
1 & $\hat{E}_{\text{low}}$  & Ile energii w~niskich częst.\ (gładkość) \\
2 & $\hat{E}_{\text{mid}}$  & Ile energii w~średnich częst.\ (tekstury) \\
3 & $\hat{E}_{\text{high}}$ & Ile energii w~wysokich częst.\ (krawędzie) \\
4 & $\hat{H}_{\text{spec}}$ & Entropia spektralna --- jak ,,równomiernie'' rozkłada się energia (wyższa = bardziej złożony obraz) \\
5 & $f_c$                   & Centroid częstotliwościowy --- ,,środek ciężkości'' widma (wyższy = więcej detali) \\
6 & $E_{\text{H}}$          & Energia krawędzi poziomych \\
7 & $E_{\text{V}}$          & Energia krawędzi pionowych \\
8 & $E_{\text{D1}}$         & Energia krawędzi diagonalnych (45\textdegree) \\
9 & $E_{\text{D2}}$         & Energia krawędzi diagonalnych ($-$45\textdegree) \\
\bottomrule
\end{tabular}
\end{table}

Z~tych 9~cech najważniejszą metryką jest \textbf{stosunek wysokich do niskich
częstotliwości} $R_{\text{HL}} = \hat{E}_{\text{high}} /
\hat{E}_{\text{low}}$:
\begin{itemize}
    \item $R_{\text{HL}} \approx 5$ --- obraz gładki, rozmyty, mało krawędzi.
    \item $R_{\text{HL}} \approx 29$ --- obraz ostry, pełen detali i~krawędzi
          narzędzi.
\end{itemize}

W~naszym zbiorze danych $R_{\text{HL}}$ waha się od~4{,}79 do~28{,}79
(zakres $6\times$!), co potwierdza ogromną różnorodność spektralną.


% -----------------------------------------------------------------------------
\subsection{Krok 4: Pipeline selekcji --- 4 ,,oczy'' patrzące na dane}
\label{sec:krok4}

Wybieramy 20\,000 obrazów nie jednym kryterium, lecz czterema
komplementarnymi:

\begin{table}[h]
\centering
\caption{Cztery grupy cech w~pipeline selekcji i~ich rola.}
\label{tab:pipeline_simple}
\begin{tabular}{lccl}
\toprule
\textbf{Cecha} & \textbf{Wymiary} & \textbf{Waga} & \textbf{Co ,,widzi''} \\
\midrule
DINO (PCA)  & 32  & 35\% & Semantyka --- \emph{co} jest na obrazie \\
EL2N        & 1   & 30\% & Trudność --- jak \emph{trudny} jest obraz dla DETR \\
SAM         & 1   & 20\% & Złożoność --- ile \emph{obiektów/segmentów} jest na obrazie \\
\textbf{Fourier} & \textbf{9} & \textbf{15\%} & Częstotliwości --- \emph{jak wygląda} obraz w~sensie ostrości/tekstury \\
\bottomrule
\end{tabular}
\end{table}

\paragraph{Dlaczego nie wystarczy samo DINO?}
DINO (sieć neuronowa typu Vision Transformer) świetnie rozumie \emph{co}
jest na obrazie (,,narzędzie chirurgiczne w~lewym górnym rogu''), ale nie
rozróżnia obrazu ostrego od rozmytego --- oba wyglądają tak samo semantycznie.
Fourier uzupełnia tę lukę.

\paragraph{Jak to działa razem?}
Wszystkie 4~wektory cech (po przeskalowaniu wagami) łączymy w~jeden wektor
na obraz.  Potem uruchamiamy \textbf{klasteryzację $k$-means} z~$k = 20\,000$
klastrów.  Z~każdego klastra wybieramy 1~obraz najbliższy centroidowi.
Wynik: 20\,000 maksymalnie różnorodnych obrazów.


% -----------------------------------------------------------------------------
\subsection{Krok 5: Dlaczego Fourier jest ważniejszy niż się wydaje?}
\label{sec:krok5}

Na pierwszy rzut oka: Fourier ma 15\% wagi, DINO ma 35\%.  Ale to nie
cała historia.

\subsubsection{Ważność per-wymiar}

Wytrenowaliśmy regresor Random Forest, który na podstawie cech DINO,
Fourier i~SAM przewiduje trudność EL2N obrazu.  Wyniki:

\begin{table}[h]
\centering
\caption{Ważność cech w~przewidywaniu trudności EL2N.}
\label{tab:importance_simple}
\begin{tabular}{lccc}
\toprule
\textbf{Cecha} & \textbf{Wymiary} & \textbf{Łączna ważność} &
    \textbf{Ważność na 1 wymiar} \\
\midrule
DINO (PCA)       & 32 & 48{,}4\% & 1{,}51\% \\
\textbf{Fourier} &  9 & \textbf{46{,}3\%} & \textbf{5{,}14\%} \\
SAM              &  1 &  5{,}3\% & 5{,}30\% \\
\bottomrule
\end{tabular}
\end{table}

\textbf{Kluczowy wniosek:} Każdy z~9 wymiarów Fouriera niesie
$5{,}14 / 1{,}51 = 3{,}4\times$ \textbf{więcej informacji} niż każdy z~32
wymiarów DINO.  Fourier jest ,,gęsty informacyjnie'' --- mało wymiarów, ale
każdy bardzo ,,mówi'' modelowi coś nowego.

\subsubsection{Dlaczego tak jest?}

Sieci neuronowe (w~tym DETR) mają wrodzone \textbf{obciążenie spektralne}
(spectral bias): uczą się najpierw rzeczy ,,gładkich'' (niskie
częstotliwości), a~ostre detale (wysokie częstotliwości) uczą się na końcu
i~z~trudem.

Obrazy, które są ,,trudne'' dla DETR (wysoki wynik EL2N), to zazwyczaj te
z~nietypową zawartością częstotliwościową:
\begin{itemize}
    \item Ostre odbicia na mokrej rogówce (nietypowo wysokie częstotliwości).
    \item Rozmyte klatki z~poruszenia (nietypowo niskie częstotliwości).
    \item Nietypowe tekstury tkanek.
\end{itemize}

Fourier te przypadki ,,widzi'' wprost, bo mierzy właśnie rozkład
częstotliwości.  DINO tego nie koduje --- dwa obrazy mogą wyglądać
,,semantycznie tak samo'' (ten sam narzędzie w~tym samym miejscu), ale jeden
jest ostry, a~drugi rozmyty.  Dla DINO to duplikat.  Dla Fouriera --- dwa
kompletnie różne obrazy.


% -----------------------------------------------------------------------------
\subsection{Krok 6: Co mówi analiza tooltip vs background?}
\label{sec:krok6}

Porównaliśmy sygnatury Fouriera obrazów z~narzędziem (tooltip) vs bez
narzędzia (background) na próbce $n=50+50$.  Najsilniejsze różnice:

\begin{table}[h]
\centering
\caption{Top 5 cech Fouriera różniących tooltip od background
         (Cohen's $d$ --- wielkość efektu; AUC --- zdolność klasyfikacyjna).}
\label{tab:tooltip_vs_bg}
\begin{tabular}{lccl}
\toprule
\textbf{Cecha} & $d$ & AUC & \textbf{Interpretacja} \\
\midrule
$E_{\text{D1}}$ (diag.\ 45\textdegree) & +0{,}92 & 0{,}64 & Tooltipy mają silniejsze krawędzie pod kątem \\
$\hat{H}_{\text{spec}}$                & +0{,}70 & 0{,}61 & Tooltipy mają bardziej złożone widmo \\
spectral\_spread                        & $-$0{,}66 & 0{,}65 & Tooltipy mają bardziej skupione widmo \\
anisotropy                              & $-$0{,}56 & --- & Background jest bardziej izotropowy \\
psd\_slope                              & +0{,}46 & 0{,}61 & Tooltipy mają płytszy spadek energii \\
\bottomrule
\end{tabular}
\end{table}

Tooltipy (narzędzia chirurgiczne) mają:
\begin{itemize}
    \item \textbf{Więcej energii w~pasmach średnio-wysokich} (50--70\% promienia) --- ostre brzegi narzędzi.
    \item \textbf{Silniejsze krawędzie diagonalne} ($d = +0{,}92$) --- czubki narzędzi pod kątem.
    \item \textbf{Wyższą entropię spektralną} --- bardziej ,,szumowe'', złożone widmo.
\end{itemize}

To znaczy, że cechy Fouriera \textbf{naturalnie kodują informację, czy na
obrazie jest narzędzie chirurgiczne} --- bez żadnych etykiet!


% -----------------------------------------------------------------------------
\subsection{Krok 7: Czy selekcja nie psuje rozkładu danych?}
\label{sec:krok7}

Wybieramy 20\,000 z~90\,389 obrazów --- $4{,}5\times$ redukcja.  Czy nie
gubimy po drodze ważnych typów obrazów?

Sprawdzamy to \textbf{testem Kolmogorowa--Smirnowa (KS)}: porównujemy
rozkład każdej cechy Fouriera w~pełnej puli (90k) vs wybranym podzbiorze (20k).
Jeśli $p > 0{,}05$, to rozkłady są statystycznie nieodróżnialne.

\begin{table}[h]
\centering
\caption{Test KS: pełna pula (90\,389) vs wybrany podzbiór (20\,000).
         Wszystkie $p \gg 0{,}05$ --- rozkłady identyczne.}
\label{tab:ks_simple}
\begin{tabular}{lccl}
\toprule
\textbf{Cecha} & $\mu_{\text{pula}}$ & $\mu_{\text{20k}}$ & wartość $p$ \\
\midrule
$\hat{E}_{\text{low}}$  & 0{,}0525 & 0{,}0524 & 0{,}506 \\
$\hat{E}_{\text{mid}}$  & 0{,}5704 & 0{,}5706 & 0{,}820 \\
$\hat{E}_{\text{high}}$ & 0{,}3769 & 0{,}3770 & 0{,}928 \\
$\hat{H}_{\text{spec}}$ & 0{,}9894 & 0{,}9895 & 0{,}410 \\
$f_c$                   & 0{,}4415 & 0{,}4416 & 0{,}356 \\
$R_{\text{HL}}$         & 8{,}468  & 8{,}471  & 0{,}991 \\
\bottomrule
\end{tabular}
\end{table}

\textbf{Wniosek:} Pipeline selekcji zachowuje pełną różnorodność
częstotliwościową.  20\,000 obrazów ,,wygląda'' spektralnie dokładnie tak jak
90\,389 --- nie gubimy ani ostrych, ani rozmytych próbek.


% =============================================================================
\subsection{Krok 8: Twarde dane --- wyniki benchmarków}
\label{sec:krok8}

Teraz najważniejsze: \textbf{czy to naprawdę działa?}  Poniżej wyniki
z~dwóch niezależnych benchmarków.

\subsubsection{Benchmark same-distribution (ten sam zbiór testowy)}

Trenujemy DETR na oryginalnym zbiorze ($\sim$100k) i~na wybranym podzbiorze
(20k), testujemy na tym samym zbiorze testowym (Cataract, Query 81):

\begin{table}[h]
\centering
\caption{Same-distribution benchmark: DETR trenowany na 100k vs 20k
         Fourier-selected (ten sam zbiór testowy, ten sam checkpoint startowy
         epoch~170).  Źródło: \texttt{BENCHMARK\_Q81\_20251231\_013537}.}
\label{tab:bench_same}
\begin{tabular}{lcccc}
\toprule
\textbf{Model} & \textbf{Epoch} & \textbf{mAP@0.5} &
    \textbf{mAP@0.5:0.95} & \textbf{F1} \\
\midrule
YOLO (100k, najlepszy) & 70  & 64{,}5\% & 48{,}9\% & 74{,}8\% \\
YOLO (100k)            & 170 & 58{,}9\% & 44{,}5\% & 73{,}2\% \\
\midrule
DETR Q81 (100k)        & 100 & 70{,}6\% & 53{,}7\% & 77{,}2\% \\
DETR Q81 (100k)        & 160 & 69{,}2\% & 54{,}5\% & 78{,}1\% \\
\textbf{DETR Q81 (100k)} & \textbf{170} & \textbf{74{,}7\%} &
    \textbf{72{,}7\%} & \textbf{80{,}2\%} \\
\midrule
DETR 20k (finetune)    & 175 & 67{,}3\% & 56{,}5\% & 77{,}8\% \\
DETR 20k (finetune)    & 190 & 68{,}2\% & 55{,}7\% & 78{,}3\% \\
DETR 20k (finetune)    & 200 & 63{,}4\% & 51{,}6\% & 77{,}1\% \\
\bottomrule
\end{tabular}
\end{table}

\paragraph{Interpretacja.}
Na zbiorze, na którym model był trenowany (same-distribution), DETR 20k
osiąga porównywalną jakość ($\sim$68\% mAP@0.5) do pełnego modelu 100k
--- \textbf{używając $4{,}5\times$ mniej danych}.  Najlepszy epoch 170
(punkt startowy fine-tuningu) ma 74{,}7\%, ale to jest checkpoint
\emph{przed} fine-tuningiem na 20k.

\subsubsection{Benchmark cross-dataset (zupełnie inny zbiór testowy)}

To jest kluczowy test: trenujemy na naszych danych, testujemy na
\textbf{zupełnie innym zbiorze} (Roboflow Cataract Surgery), którego model
nigdy nie widział.  To sprawdza \textbf{generalizację} --- czy model nauczył
się ,,rozumieć'' narzędzia, czy tylko zapamiętał konkretne obrazy.

\begin{table}[h]
\centering
\caption{Cross-dataset benchmark (Roboflow Cataract Surgery --- model
         NIGDY nie widział tych obrazów).  Źródło:
         \texttt{BENCHMARK\_ALL\_EPOCHS\_20260112\_194440}.}
\label{tab:bench_cross}
\begin{tabular}{llcccc}
\toprule
\textbf{Model} & \textbf{Dane} & \textbf{Epoch} & \textbf{mAP@0.5} &
    \textbf{F1} & \textbf{Precision} \\
\midrule
YOLO & 100k orig.   & 80 (najl.)  & $<$0{,}1\% & 9{,}2\% & 43{,}5\% \\
YOLO & 20k finetune & 200         & 19{,}9\%   & 25{,}6\% & 71{,}9\% \\
\midrule
DETR & 100k orig.   & 20 (najl.)  & 8{,}9\%    & 23{,}2\% & 26{,}0\% \\
DETR & 100k orig.   & 170         & 2{,}0\%    & 14{,}5\% & 14{,}6\% \\
\midrule
\textbf{DETR} & \textbf{20k finetune} & \textbf{180} &
    \textbf{74{,}6\%} & \textbf{50{,}6\%} & 50{,}7\% \\
DETR & 20k finetune & 210 & 81{,}5\% & 46{,}5\% & 48{,}4\% \\
DETR & 20k finetune & 230 & 83{,}7\% & 47{,}9\% & 53{,}9\% \\
DETR & 20k finetune & 260 & \textbf{86{,}7\%} & 43{,}7\% & \textbf{59{,}3\%} \\
DETR & 20k finetune & 290 & 86{,}7\% & 44{,}2\% & 60{,}2\% \\
\bottomrule
\end{tabular}
\end{table}

\subsubsection{Co te liczby oznaczają?}

\begin{enumerate}
    \item \textbf{DETR na 100k danych prawie nie generalizuje.}
          Najlepszy wynik to 8{,}9\% mAP@0.5 na nowym zbiorze (epoch~20!).
          Im dłużej trenujemy na 100k, tym \emph{gorzej} radzi sobie na
          nowych danych (epoch~170: zaledwie 2\%).  Model się
          \textbf{przetrenowuje} (overfitting) na konkretne obrazy zamiast
          uczyć się rozpoznawać narzędzia.

    \item \textbf{Fine-tuning na 20k Fourier-selected daje skok
          +84{,}6\,pp mAP@0.5.}
          Z~2{,}0\% (epoch~170, przed fine-tuningiem) do 86{,}7\%
          (epoch~260, po fine-tuningu na 20k).  To poprawa o ponad
          $40\times$!

    \item \textbf{20k > 100k.}  Model trenowany na 20\,000 \emph{starannie
          wybranych} obrazach jest \textbf{dramatycznie lepszy} na nowych
          danych niż model trenowany na pełnych 100\,000 losowych obrazach.

    \item \textbf{YOLO na 20k też się poprawia}, ale mniej spektakularnie
          (0{,}1\% $\to$ 19{,}9\%).  DETR korzysta z~inteligentnej selekcji
          znacznie bardziej niż YOLO.
\end{enumerate}


% -----------------------------------------------------------------------------
\subsection{Krok 9: Dlaczego to działa? (Podsumowanie mechanizmu)}
\label{sec:krok9}

Trzy powody, dla których Fourier-aware selekcja tak dramatycznie poprawia
generalizację DETR:

\subsubsection{Powód 1: Eliminacja redundancji}

Wideo chirurgiczne przy 25--30 kl./s generuje masy prawie identycznych
klatek.  Fourier-based similarity filtering (próg 0{,}7) usuwa 10--20\%
duplikatów \emph{zanim} w~ogóle zaczniemy kosztowną ekstrakcję cech GPU.
Model nie marnuje czasu na uczenie się ,,tego samego'' tysiąc razy.

\subsubsection{Powód 2: Pokrycie spektralne}

Selekcja czysto semantyczna (np.\ tylko DINO) mogłaby wybrać 20\,000
obrazów, które semantycznie pokrywają cały zbiór (różne narzędzia, fazy
operacji), ale są wszystkie \textbf{w~podobnej ostrości}.  Dodanie Fouriera
wymusza, że w~wybranych 20k znajdą się zarówno:
\begin{itemize}
    \item Obrazy ostre (wysokie $R_{\text{HL}}$) --- uczą model precyzyjnej
          lokalizacji krawędzi.
    \item Obrazy rozmyte (niskie $R_{\text{HL}}$) --- uczą model radzić
          sobie z~nieidealnymi warunkami.
    \item Obrazy z~nietypowymi teksturami --- zwiększają odporność na nowe
          warunki.
\end{itemize}

\subsubsection{Powód 3: Walka z~przetrenowaniem (overfitting)}

Redundantne dane = przetrenowanie.  Sieci neuronowe ,,zapamiętują''
powtarzające się wzorce zamiast uczyć się ogólnych reguł.  100k obrazów
z~masą duplikatów powoduje, że DETR zapamiętuje konkretne klatki
(stąd: 2\% na nowym zbiorze).  20k różnorodnych obrazów zmusza model do
uczenia się \textbf{cech ogólnych} (stąd: 86{,}7\% na nowym zbiorze).


% =============================================================================
\subsection{Krok 10: Podsumowanie w~jednym zdaniu}
\label{sec:krok10}

\bigskip
\noindent\rule{\textwidth}{0.4pt}

\noindent\textbf{Podsumowanie:}
Transformata Fouriera pozwala zmierzyc ,,jak wyglada'' obraz
w~sensie ostrosci, tekstury i~krawedzi.  Uzywajac tej informacji do
inteligentnego wyboru 20\,000 z~90\,000 obrazow treningowych, osiagamy
DETR, ktory generalizuje $40\times$ lepiej na nowych, niewidzianych danych
niz model trenowany na wszystkich 90\,000 obrazach.

\noindent\rule{\textwidth}{0.4pt}
\bigskip

\newpage
